\chapter{Generalizations and additional considerations}
\label{cha:generalized-spectral}

So far, we have seen that the spectral method, along with appropriate statistical models, provides a simple yet effective workhorse for statistical estimation. However, the vanilla form of the spectral method, which directly extracts the top eigenspace or singular space of certain data matrices, may not be sufficient. In this chapter, we illustrate modifications to the spectral method, which are often found necessary to further enhance sample efficiency, increase robustness to outliers, incorporate signal priors, and deal with asymmetric perturbations.

%
\section{Truncated spectral method for sample efficiency} \label{sec:improved-spectral-truncated}

The generic recipe for spectral methods described in the previous chapters works well when one has  sufficient samples compared to the underlying signal dimension. It might not be effective though if the sample complexity is on the order of  the information-theoretic limit. 

In what follows, we use phase retrieval, discussed in Section~\ref{sec:phase_retrieval}, to demonstrate the underlying issues and how to address them\footnote{Strategies of similar spirit have been proposed for other problems; see, e.g.,~\citep{keshavan2010matrix} for matrix completion. }. Recall that Theorem~\ref{thm:PR-L2-performance} requires a sample complexity $m \gtrsim n  \log^3 n$, which is a few logarithmic factor larger than the  signal dimension $n$.  
%For matrix completion, we assume the sampling probability $p > \log^3 n/ n$ in Theorem~\ref{thm:spectral_MC}.
%
What happens if we only have access to $m \asymp n$ samples, which is the information-theoretic limit (order-wise) for phase retrieval? In this more challenging regime, it turns out that we have to modify the standard recipe by applying appropriate preprocessing before forming the surrogate matrix $\bm{M}$ in \eqref{eq:D_PR_org}.


We start by explaining why the surrogate matrix $\bm{M}$ in \eqref{eq:D_PR_org} for phase retrieval must be suboptimal in terms of sample complexity. For all $1 \le j \le m$,
\begin{equation*}
	\norm{\bm{M}} \ge \frac{\va_j^\top \bm{M} \va_j}{\va_j^\top \va_j}=  \frac{1}{m} \sum_{i=1}^m y_i \frac{(\va_i^\top \va_j)^2}{\va_j^\top \va_j} \geq \frac{1}{m}  y_j  \| \va_j \|_2^2.
\end{equation*}
%
In particular, taking $j = i^\ast  = \underset{{1 \le i \le m}}{\arg \max} \ y_i$ gives us
\begin{equation}\label{eq:bound_id2}
	\norm{\bm{M}} \ge \frac{(\max_i y_i) \| \va_{i^\ast} \|_2^2}{m}.
\end{equation}
Under the Gaussian design, $\set{y_i / \| \vxp \|_2^2}$ is a collection of i.i.d.~$\chi^2$ random variables with 1 degree of freedom. It follows from well-known estimates in extreme value theory \citep{Ferguson:1996} that 
\[
\max_{1\le i \le m} \, y_i \approx 2\|\vxp\|_2^2 \log m,\qquad \text{as}\; m \to \infty.
\] 
Meanwhile, note that $\| \va_{i^\ast} \|_2^2  \approx n$ for $n$ sufficiently large. It follows from \eqref{eq:bound_id2} that
\begin{equation}\label{eq:PR_sc}
	\norm{\bm{M}} \ge \big(1+o(1)\big)  \frac{ n \log m}{  m} \|\vxp\|_2^2  .
\end{equation}
%

%These estimates can be made precise as follows.
%\begin{lemma}\label{lemma:T_id}
%Suppose the sensing vectors $\set{\va_i}$ are drawn from the i.i.d. Gaussian distribution. Then
%\[
%\norm{\mD - \mathbb{E}[\mD]} > \norm{\vxp}^2 (n \log m) / m - 2 \norm{\vxp}^2  - 1
%\]
%with probability at least $1 - \mathcal{O}(1/n^2)$.
%\end{lemma}

%Recall that $\mathbb{E}[\bm{Y}] = \vxp \vxp^\top + \| \vxp\|_2^2 \boldsymbol{I}_n$.  Its spectral norm is bounded and its leading eigenvector is exactly the true object $\vxp$. Our goal is to make sure the noisy matrix $\bm{Y}$ is close to $\EE[\bm{Y}]$ so that the principal eigenvector of $\bm{Y}$ stays close to that of $\EE[\bm{Y}]$. 
% = \vxp \vxp^\top + \| \vxp\|_2^2 \boldsymbol{I}_n

Recall that $\mathbb{E}[\bm{M}] = 2 \vxp \vxp^\top +  \| \vxp\|_2^2 \boldsymbol{I}_n$ has a bounded spectral norm, then \eqref{eq:PR_sc} implies that, to keep the deviation between $\bm{M}$ and $\mathbb{E}[\bm{M}]$ well-controlled, we must at least have
\[
\frac{n \log m}{ m} \lesssim 1.
\]
This condition, however, cannot be satisfied when we have linear sample complexity $m \asymp  n$. This explains why we need a sample complexity $m \gg n \log n$ in Theorem~\ref{thm:PR-L2-performance}. 
%Indeed, in \cite{netrapalli2013phaseAM}, the authors established performance guarantees for the spectral method by assuming $m =  c \cdot n \log^3 n$ (for some sufficiently large $c$), and this requirement was later reduced to $m = c \cdot n \log n$ in \cite{candes2015phase}. 

The above diagnosis also suggests an easy fix: since the main culprit lies in the fact that $\max_i y_i$ is unbounded (as $m \to \infty$), we can apply a preprocessing function $\mathcal{T}(\cdot)$ to $y_i$ to keep the quantity bounded. Indeed, this is the key idea behind the {\em truncated spectral method} proposed by \cite{chen2015solving}, in which the surrogate matrix is modified as
%
\begin{equation}
	\label{eq:D_PR_T}
	\bm{M}_{\mathcal{T}} := \frac{1}{m} \sum_{i=1}^m \mathcal{T}(y_i) \va_i \va_i^\top,
\end{equation} 
%
where 
\begin{equation}\label{eq:trimming}
\mathcal{T}(y) := y \charfn_{\set{\abs{y} \le  \frac{\gamma}{m} \sum_{i=1}^m y_i}}
\end{equation}
for some predetermined truncation threshold $\gamma$. The initial point $\bm{x}_0$ is then formed by scaling the leading eigenvector of $\bm{M}_{\mathcal{T}}$ to have roughly the same norm of $\bx_{\star}$, which can be estimated by $\overline{y} =\frac{1}{m} \sum_{i=1}^m y_i $. This is essentially performing a trimming operation, removing any entry of $y_i$ that bears too much influence on the leading eigenvector.  The trimming step turns out to be very effective, allowing one to achieve order-wise optimal sample complexity. The following theorem is due to \cite{chen2015solving}. 
%
%\begin{align}
%	\label{eq:init-spectral-proc}
%	\bm{x}_0 = \sqrt{\lambda}\bm{u},
%\end{align}
%
%where $\lambda$ and $\bm{u}$ are the leading eigenvalue and eigenvector of $\bm{Y}_{\mathcal{T}}$, respectively. 


%
\begin{theorem}[$\mathsf{Truncated~spectral~method~for~phase~retrieval}$]
	\label{thm:truncated-PR}
	Consider phase retrieval in Section~\ref{sec:phase_retrieval}. Fix any $\zeta > 0$, and suppose $m \ge c_0 n $ for some sufficiently large constant $c_0>0$.  Then the truncated spectral estimate obeys
\[
	\mathsf{dist}(\bm{x},\bm{x}^{\star})  \le \zeta \| \vxp \|_2
\]
with probability at least $1 - O(n^{-2})$.
\end{theorem}

Subsequently, several different designs of the preprocessing function have been proposed in the literature. One example was given in \cite{wang2017solving}, where
\begin{equation}\label{eq:subset}
	\mathcal{T}(y)  = \charfn_{\set{y \ge \gamma  }},
\end{equation}
%
where $\gamma$ is the $(cm)$-largest value (e.g.~$c=1/6$) in $\{ y_j \}_{1\le j\le m}$. In words, this method only employs a subset of design vectors that are better aligned with the truth $\bm{x}_{\star}$.  By properly tuning the parameter $c$, this truncation scheme performs competitively as the scheme in \eqref{eq:trimming}.

\begin{figure}[ht]
\begin{center}
\includegraphics[width=0.8\textwidth]{figures/stanford} \\
(a) ground truth \\
\includegraphics[width=0.8\textwidth]{figures/stanford_quad_init_WF} \\
(b) initialization via the vanilla spectral method \\
\includegraphics[width=0.8\textwidth]{figures/stanford_quad_init_TWF}  \\
(c) initialization via the truncated spectral method
\end{center}
	\caption{Illustration of signal recovery using the vanilla spectral method and the truncated spectral method in phase retrieval. Here, \textbf{add parameter settings}. } \label{fig:phase_retrieval_twf}
\end{figure}



%
\section{Truncated spectral method for removing sparse outliers}

When the samples are susceptible to adversarial entries, the spectral method might not work properly even with the presence of a single outlier whose magnitude can be arbitrarily large to perturb the leading eigenvector of the data matrix $\bm{M}$. Therefore, the spectral method needs to be modified carefully in order to still function properly in the presence of outliers.

\paragraph{Median-truncated spectral method.} One strategy to remove outliers is to truncate away spurious samples based on the sample median, as the median is known to be robust against arbitrary outliers \citep{zhang2016provable,li2017nonconvex}. We illustrate this median-truncation approach through an example of robust phase retrieval \citep{zhang2016provable}, where we assume a subset of samples in \eqref{eq:PR-samples} is corrupted arbitrarily, with their index set denoted by $\mathcal{S}$ with $|\mathcal{S}|=\alpha  m$, $\alpha \in (0,1/2]$ is the faction of outliers. Mathematically, the measurement  model in the presence of outliers is given by
%
\begin{equation} \label{eq:robust_pr_model}
	y_i = \begin{cases}
(\ba_i^\top\bx_{\star})^2, \quad  & i \notin \mathcal{S}, \\
\mathsf{arbitrary} , & i \in \mathcal{S}. 
	\end{cases}
\end{equation}
%
The goal is to still recover $\bm{x}_{\star}$ in the presence of many outliers (e.g.~a constant fraction of measurements are outliers). To mitigate this issue, a median-truncation scheme was proposed in \cite{zhang2016provable,li2017nonconvex}, where
\begin{equation}\label{eq:T_median}
	\mathcal{T}(y)  = y \charfn_{\set{y_i \leq \gamma \, \mathsf{median}\{ y_j \}_{j=1}^m}}
\end{equation}
for some predetermined constant $\gamma >0$, $\mathsf{median}\{ y_j \}_{j=1}^m $ is the median of the samples $\{ y_j \}_{j=1}^m$. The data matrix $\bM$ is then constructed via \eqref{eq:D_PR_T} using the above preprocessing function.

 By including only a subset of samples whose values are not excessively large compared with the sample median of the samples, the preprocessing function in \eqref{eq:T_median} makes the spectral method more robust against sparse and large outliers.
The following theorem is due to \cite{zhang2016provable}. 
\begin{theorem} 
	\label{thm:truncated-PR-init}
	Consider the robust phase retrieval problem in \eqref{eq:robust_pr_model}, and fix any $\zeta > 0$. There exist some constants  $c_0, c_1>0$ such that if 
$m\geq c_{0}n$ and $\alpha\leq c_1$, then the median-truncated spectral estimate obeys
\[
	\mathsf{dist}(\bm{x},\bm{x}^{\star})  \le \zeta \| \vxp \|_2
\]
with probability at least $1 - O(n^{-2})$.
\end{theorem}

Theorem~\ref{thm:truncated-PR-init} suggests that even with a constant fraction of outliers, the median-truncated spectral method can still robustly estimate $\bm{x}^{\star}$ with a near-optimal number of measurements $m=O(n)$, thus further enhancing the resilience of the truncated spectral method in the presence of adversarial corruptions.

%median-TWF needs O(n) sample for initialization.
\paragraph{Hard-thresholded spectral method.}
The idea of applying truncation to form a spectral estimate is also used in the robust PCA problem \citep{chandrasekaran2011siam,candes2009robustPCA,vaswani2017robust}, when we are provided an upper bound on the fraction of outliers. 
Suppose we are given a matrix $\bm{\Gamma}_{\star}\in\mathbb{R}^{n_1\times n_2}$ that is a superposition of a rank-$r$ matrix $\bm{M}_{\star}\in\mathbb{R}^{n_1\times n_2}$ and a sparse matrix $\bm{S}_{\star}\in\mathbb{R}^{n_1\times n_2}$:
%
\begin{equation}\label{eq:rpca_decomp}
\bm{\Gamma}_{\star} = \bm{M}_{\star} + \bm{S}_{\star}.
\end{equation}
%
The goal is to separate $\bm{M}_{\star}$ and $\bm{S}_{\star}$ from the (possibly partial) entries of $\bm{\Gamma}_{\star}$. This problem is also known as {\em robust PCA}, since we can think of it as recovering the low-rank factors of  $\bm{M}_{\star}$ when the observed entries are corrupted by {\em sparse outliers} (modeled by $\bm{S}_{\star}$). The problem spans numerous applications in computer vision, medical imaging, and video surveillance. 

Similar to the matrix completion problem, we assume the random sampling model \eqref{assump:mc-bernoulli}, where $\Omega$ is the set of observed entries. 
In order to make the problem well-posed, we need the incoherence parameter of $\bM_{\star}$ as defined in Definition~\ref{assump:mc-incoherence} as well, which precludes $\bM_{\star}$ from being too spiky.  In addition, it is sometimes convenient to introduce the following deterministic condition on the sparsity pattern and sparsity level of $\bm{S}_{\star}$,  originally proposed by \cite{chandrasekaran2011siam}. Specifically, it is assumed that the non-zero entries of $\bm{S}_{\star}$ are ``spread out'', where there are at most a  fraction $\alpha$ of non-zeros per row\,/\,column. Mathematically, this means:

\begin{assumption}
	It is assumed that $\bS_{\star}\in\mathcal{S}_\alpha\subseteq \mathbb{R}^{n_1\times n_2}$, where
%
\begin{equation}
	\mathcal{S}_\alpha := \left\{\bS \,:\,  \|(\bS_{\star})_{i,\cdot}\|_0 \leq \alpha n_2, \; \|(\bS_{\star})_{\cdot,j}\|_0 \leq \alpha n_1, \forall i,j \right\}. 
\end{equation}
%
\end{assumption}


Since the observations are also potentially corrupted by large but sparse outliers, it is useful to first clean up the observations $(\bm{\Gamma}_{\star})_{i,j} = (\bm{M}_{\star})_{i,j} + (\bm{S}_{\star})_{i,j}$ before constructing the matrix $\bM$ as in \eqref{eq:rescaled-M-matrix-completion-l2}. Indeed, this is the strategy proposed in \cite{yi2016fast}. 
%For simplicity, we present in what follows only the results for the fully observed case. 
To proceed, we start by forming an estimate of the sparse outliers via the hard-thresholding operation as
\begin{equation}\label{eq:hard_threshodling_init}
\bm{S}_{0}=\mathcal{H}_{c\alpha np}\big(\mathcal{P}_{\Omega}(\bm{\Gamma}_{\star} )\big).
\end{equation}
%We start by forming an estimate of the sparse outliers. Let
%\[
%\mathcal{T}_{\alpha}[\bm{\Gamma}] = \begin{cases}
%\bm{\Gamma}_{i,j}, &\text{if } \abs{\bm{\Gamma}_{i,j}} \ge \abs{\bm{\Gamma}_{(i, \cdot)}^{(\alpha n_2)}} \text{ and } \abs{\bm{\Gamma}_{i,j}} \ge \abs{\bm{\Gamma}_{(\cdot, j)}^{(\alpha n_1)}} \\
%0,&\text{otherwise}
%\end{cases},
%\]
%where $\bm{\Gamma}_{(i, \cdot)}^{(k)}$ (resp. $\bm{\Gamma}_{(\cdot, j)}^{(\alpha n_1)}$) denotes the $k$th largest entry of $\bm{\Gamma}_{(i, \cdot)}$ (resp. $\bm{\Gamma}_{(\cdot, j)}$). 
where $c>0$ is some predetermined constant (e.g.~$c=3$) and $\mathcal{P}_{\Omega}$ is defined in \eqref{defn:Pomega}. Here, $\mathcal{H}_l(\cdot)$ is the hard-thresholding operator defined as
%
\[
	[\mathcal{H}_{l}(\bA)]_{j,k} := \begin{cases}  A_{j,k},  &\text{if }|A_{j,k}|\geq |A_{j,\cdot}^{(l)}| \text{ and }
							|A_{j,k}|\geq |A_{\cdot,k}^{(l)}|,  \\
							0, & \text{otherwise,}
					  \end{cases}
\]
where 
$A_{j,\cdot}^{(l)}$ (resp.~$A_{\cdot,k}^{(l)}$) denotes the $l$th largest entry (in magnitude) in the $j$th row (resp.~column) of $\bm{A}$.
The  idea is simple: an entry is likely to be an outlier if it is simultaneously among the largest entries in the corresponding row and column. Armed with this estimate, we form the matrix $\bM$ as
%
\begin{equation}
	\label{eq:Y-RPCA}
	\bm{M} =\frac{1}{p} \mathcal{P}_{\Omega}(\bm{\Gamma}_{\star}   - \bm{S}_{0}).
\end{equation}

 
%	


One can then apply the spectral method to $\bm{M}$ in \eqref{eq:Y-RPCA} (similar to the matrix completion case). This approach enjoys the following performance guarantee, due to \cite{yi2016fast}; see also \citep{cherapanamjeri2017nearly}.
%
\begin{theorem} \label{thm:spectral_RPCA}
Suppose that the robust PCA problem. If the sample size and the sparsity fraction satisfy $n_1  p\geq c_0  \kappa^2 \mu r  \log n_2$ and $\alpha\leq c_1/( \mu \kappa  r )$ for
	some sufficiently large (resp. small) constant $c_0>0$ (resp. $c_1>0$), then 
with probability at least $1-O\left(n^{-1}\right)$,  
%
	\[
  \max\Big\{ \mathsf{dist}\left(\bm{U},\bm{U}^{\star}\right),\mathsf{dist}\left(\bm{V},\bm{V}^{\star}\right) \Big\}  \leq  c_2 \alpha \mu r \kappa + c_3 \kappa \sqrt{\frac{\mu r \log n_2}{n_1 p}},
	\]	
for some constants $c_2, c_3>0$.
\end{theorem}

It can be seen that Theorem~\ref{thm:spectral_RPCA} reduces to Theorem~\ref{thm:mc-l2-subspace} for the case of matrix completion when there are no outliers, i.e. $\alpha =0$. Therefore, the truncated spectral method injects additional robustness properties to the spectral method without hurting its performance in the outlier-free case, hence it is resilient in an oblivious fashion.   




%
\section{Truncated spectral method for sparse phase retrieval}\label{sec:spectral_init_spr}

Next, we briefly discuss how the spectral method can be modified to incorporate structural priors. In many problems of practical interest, we might be given some prior knowledge about the signal of interest, encoded by a constraint set $\bm{x}_{\star}\in\mathcal{C}\subset \mathbb{R}^n$. Therefore, it is natural to apply projection to enforce the desired structural constraints. A special structure that will receive special attention is the sparsity prior, where it is known {\em a priori} that the signal of interest $\bm{x}_{\star} $ is $k$-sparse, where $k\ll n$. We shall illustrate this idea through the application of sparse phase retrieval.


\paragraph{Sparse phase retrieval.} Consider the phase retrieval problem formulated in Section~\ref{sec:formulation-PR}, where we now additionally assume $\bx_{\star}$ is  $k$-sparse, i.e. $\| \bx_{\star}\|_0 = k$.  The sparse phase retrieval problem has numerous applications in computational imaging \citep{soltanolkotabi2017structured,shechtman2014gespar,chen2015exact}, as it offers the premise to further reducing the sampling cost by exploiting the sparsity structure.
A simple idea is to first identify the support of $\bx_{\star}$, and then try to estimate  the nonzero values by applying the spectral method over the submatrix of $\bM$ on the estimated support. Towards this end, we recall that 
$$\mathbb{E}[\bM] = 2\bx_{\star}\bx_{\star}^{\top}+ \|\bx_{\star}\|_2^2\bI_n, $$ 
and therefore the larger ones of the diagonal entries of $\bM$ are more likely to be in the support of $\bx_{\star}$. In light of this, a simple thresholding strategy adopted in \cite{cai2016optimal} is to compare $M_{i,i}$ against some preset threshold $\gamma$:
$$ \hat{\mathcal{S}} = \{i: \; Y_{i,i}> \gamma \} .$$
The nonzero part of $\bx_{\star}$ is then found by applying the spectral method outlined in Section~\ref{sec:phase_retrieval} to 
$$\bm{M}_{\hat{\mathcal{S}}} =\frac{1}{m}\sum_{i=1}^m y_i \bm{a}_{i,\hat{\mathcal{S}}}\bm{a}_{i,\hat{\mathcal{S}}}^{\top}, $$ 
where $\bm{a}_{i,\hat{\mathcal{S}}}$ is the subvector of $\bm{a}_{i}$ coming from the support $\hat{\mathcal{S}}$. 
In short, this strategy provably leads to a reasonably good initial estimate, as long as the sample size $m \gtrsim k^2 \log m$ is sufficiently large; the exact form of the theory can be found in \cite{cai2016optimal}. 
See also \citep{wang66sparse} for a more involved approach, which provides better empirical performance.
 
 
\begin{remark}
Careful readers might realize that the sample complexity of the sparse spectral method is much larger than the information-theoretical lower bound established in, e.g., \cite{eldar2014phase}. 
Unfortunately, to date there is no tractable  procedure that can provably guarantee a consistent estimate of $\bm{x}_\star$ when $m = O( k\log n) $. Rather, all known computationally feasible algorithms (both convex and nonconvex)  require a sample complexity at least on the order of $k^2 \log n$ under the Gaussian design (i.e. Assumption~\ref{assump:pr-gaussian}), unless $k$ is sufficiently large or other structural information is available \citep{li2013sparse,oymak2012simultaneously,chen2015exact,cai2016optimal,jagatap2017fast}.  All in all, there seems to exist a computational barrier for sparse phase retrieval between the sample complexity requirements for computationally efficient procedures and the information-theoretic limit. 
\end{remark}

 

\input{chapters/asymmetric.tex}

\section{Notes}

The first order-wise optimal spectral method  for phase retrieval was proposed by  \citet{chen2015solving}, based on the truncation idea. This method has multiple variants \citep{li2017nonconvex, wang2017solving}, and has been shown to be robust against noise. The precise asymptotic characterization of the spectral method was first obtained in \citet{lu2020phase}, where its sensitivity to model mismatch is studied in \citet{monardo2019sensitivity}. Based on this characterization, \citep{mondelli2017fundamental,luo2019optimal} later devised optimal designs of spectral methods in phase retrieval when the sensing matrix follows the Gaussian design. \cite{ma2019spectral,dudeja2020analysis} considered similar questions when the sensing matrix
is Haar distributed (e.g. an isotropically random unitary matrix).

\citep{fannjiang2020numerics} provided an extensive discussion on initialization strategies for algorithmic phase retrieval, including but not limited to various forms of the spectral method.

Sparse phase retrieval \citep{jagatap2019sample} 

sparse phase retrieval \citep{ma2019spectral}
 
\cite{yuan2013truncated}

sparse spectral method for blind calibration \cite{li2018blind}

Lee and Wu \cite{lee2017near}


 ranking from pairwise comparisons \citep{negahban2016rank,chen2015spectral,chen2017spectral}, 

sparse PCA \cite{amini2008high,JohLu09}